% References, IEEE style, numbered in alphabetical order of first author.
% Each entry is copied from the paper itself.
% Numbering follows the order below.

\begin{thebibliography}{99}
\addcontentsline{toc}{chapter}{References}

\bibitem{bertaux2016}
J.-L. Bertaux, I. V. Khatuntsev, A. Hauchecorne, W. J. Markiewicz, E. Marcq,
S. Lebonnois, M. Patsaeva, A. Turin, and A. Fedorova, ``Influence of Venus topography
on the zonal wind and UV albedo at cloud top level: the role of stationary gravity
waves,'' \emph{Journal of Geophysical Research: Planets}, vol. 121, pp. 1087 to 1101,
2016, doi: 10.1002/2015JE004958.

\bibitem{bierson2020}
C. J. Bierson and X. Zhang, ``Chemical cycling in the Venusian atmosphere: a full
photochemical model from the surface to 110 km,'' \emph{Journal of Geophysical Research:
Planets}, vol. 125, art. e2019JE006159, 2020, doi: 10.1029/2019JE006159.

\bibitem{dai2023kzz}
L. Dai, W. Shao, H. Gu, and Z. Sheng, ``Determination of the eddy diffusion in the
Venusian clouds from VeRa sulfuric acid observations,'' \emph{Astronomy and
Astrophysics}, vol. 679, art. A155, 2023, doi: 10.1051/0004-6361/202347714.

\bibitem{dai2024}
L. Dai, W. Shao, and Z. Sheng, ``An investigation into Venusian atmospheric chemistry
based on an open-access photochemistry-transport model at 0 to 112 km,''
\emph{Astronomy and Astrophysics}, vol. 689, art. A55, 2024,
doi: 10.1051/0004-6361/202450552.

\bibitem{dai2025review}
L. Dai, D. V. Titov, W. D. Shao, X. Zhang, J. Cui, and S. Fan, ``Venus cloud research:
progress and perspectives,'' \emph{Space Science Reviews}, vol. 221, art. 51, 2025,
doi: 10.1007/s11214-025-01176-4; preprint arXiv:2506.06164.

\bibitem{egan2025osso}
J. V. Egan, W. Feng, A. D. James, J. Manners, D. R. Marsh, S. Lebonnois, F. Lef\`evre,
A. Stolzenbach, and J. M. C. Plane, ``Is OSSO a significant contributor to the unknown
UV absorber in Venus' atmosphere?,'' \emph{Geophysical Research Letters}, vol. 52,
art. e2024GL113090, 2025, doi: 10.1029/2024GL113090.

\bibitem{egan2025lab}
J. V. Egan, A. D. James, and J. M. C. Plane, ``Laboratory measurements of ferric
chloride (FeCl\textsubscript{3}) under Venusian conditions,'' \emph{ACS Earth and Space
Chemistry}, 2025, doi: 10.1021/acsearthspacechem.5c00132.

\bibitem{egan2026}
J. V. Egan, W. Feng, A. D. James, J. Manners, D. R. Marsh, S. Lebonnois, F. Lef\`evre,
A. Stolzenbach, and J. M. C. Plane, ``Ferric chloride: the Venusian unknown UV
absorber?,'' preprint, ESS Open Archive, posted 14 January 2026, not peer reviewed,
doi: 10.22541/essoar.176840811.14859582/v1.

\bibitem{jiang2024}
C. Z. Jiang, P. B. Rimmer, G. G. Lozano \emph{et al.}, ``Iron-sulfur chemistry can
explain the ultraviolet absorber in the clouds of Venus,'' \emph{Science Advances},
vol. 10, art. eadg8826, 2024, doi: 10.1126/sciadv.adg8826.

\bibitem{lefevre2022}
M. Lef\`evre, E. Marcq, and F. Lef\`evre, ``The impact of turbulent vertical mixing in
the Venus clouds on chemical tracers,'' \emph{Icarus}, vol. 386, art. 115148, 2022,
doi: 10.1016/j.icarus.2022.115148.

\bibitem{lefevre2024}
M. Lef\`evre, F. Lef\`evre, E. Marcq, A. M\"a\"att\"anen, A. Stolzenbach, and N. Streel,
``Impact of the turbulent vertical mixing on chemical and cloud species in the Venus
cloud layer,'' \emph{Geophysical Research Letters}, vol. 51, art. e2024GL108771, 2024,
doi: 10.1029/2024GL108771.

\bibitem{limaye2018}
S. S. Limaye, R. Mogul, D. J. Smith, A. H. Ansari, G. P. S\l{}owik, and
P. Vaishampayan, ``Venus' spectral signatures and the potential for life in the
clouds,'' \emph{Astrobiology}, vol. 18, no. 9, pp. 1181 to 1198, 2018,
doi: 10.1089/ast.2017.1783.

\bibitem{perezhoyos2018}
S. P\'erez-Hoyos, A. S\'anchez-Lavega, A. Garc\'ia-Mu\~noz, P. G. J. Irwin, J. Peralta,
G. Holsclaw, W. M. McClintock, and J. F. Sanz-Requena, ``Venus upper clouds and the UV
absorber from MESSENGER/MASCS observations,'' \emph{Journal of Geophysical Research:
Planets}, vol. 123, pp. 145 to 162, 2018, doi: 10.1002/2017JE005406.

\bibitem{rimmer2021}
P. B. Rimmer, S. Jordan, T. Constantinou, P. Woitke, O. Shorttle, R. Hobbs, and
A. Paschodimas, ``Hydroxide salts in the clouds of Venus: their effect on the sulfur
cycle and cloud droplet pH,'' \emph{The Planetary Science Journal}, vol. 2, art. 133,
2021, doi: 10.3847/PSJ/ac0156.

\bibitem{spacek2026}
J. Spacek, P. B. Rimmer, J. J. Petkowski, and Y. J. Lee, ``A model of UV-blue
absorbance in bulk liquid of Venusian cloud aerosols is consistent with efficient
organic absorbers at high concentrations,'' \emph{Astrobiology}, vol. 26, no. 9,
pp. 727 to 739, 2026, doi: 10.1177/15311074261477502.

\bibitem{titov2008}
D. V. Titov, F. W. Taylor, H. Svedhem, N. I. Ignatiev, W. J. Markiewicz, G. Piccioni,
and P. Drossart, ``Atmospheric structure and dynamics as the cause of ultraviolet
markings in the clouds of Venus,'' \emph{Nature}, vol. 456, pp. 620 to 623, 2008,
doi: 10.1038/nature07466.

\bibitem{trabelsi2026}
T. Trabelsi and J. S. Francisco, ``Chlorine-sulfur isomers as parents of
ClS\textsubscript{2} and SCl\textsubscript{2} on Venus: spectroscopy and photochemistry
of ClSSCl, SSCl\textsubscript{2}, and (ClS)\textsubscript{2},'' \emph{The Journal of
Chemical Physics}, vol. 164, art. 084302, 2026, doi: 10.1063/5.0317311.

\end{thebibliography}
